\subsection{The Call-By-Value $\lambda$-calculus}
\label{ss:cbv}

The grammar
of the untyped call-by-value $\lambda$-calculus, $\LasV$, 
has values $V$, terms $M$, evaluation contexts $E$, and general contexts $C$:
\[
\begin{array}{llll}
\Values &  V & \defeq & x \midd \lambda x.M \\
%   
\Terms & M,N & \defeq & V \midd M N \\
% 
\EContexts & E & \defeq & \contexthole \midd V E \midd E M \\
%
\Contexts & C & \defeq & \contexthole \midd \lambda x.C \midd M C \midd C M
\end{array}\]
 where $\contexthole$ stands for the hole of a context. 
The call-by-value reduction $\redv$ has two rules:
  \begin{mathpar}
   \inferrule*{ }{(\lambda x.M) V \redv M\subst{x}{V}} \qquad
   \inferrule*{M \redv N}{E[M] \redv E[N]}
 \end{mathpar}
 
 In the following, we write %$M\Diverge$ to indicate that a term diverge, and
 $M \Converge M'$ to indicate that $M \Redv M'$ with $M'$ an eager normal form,
 that is, either a value or a stuck call $E[x V]$.

 We recall  Lassen's 
\emph{eager normal-form (enf) bisimilarity} \cite{lassentrees}.
Enf-bisimulations relate terms, values, and evaluation contexts.

\begin{definition}
\label{d:enf}
 An \emph{enf-bisimulation} is a triple of relations on terms $\RR_{\enfM}$,
 values $\RR_{\enfV}$, and evaluation contexts $\RR_{\enfK}$ that satisfies:
 
 \begin{itemize}
  \item $M_1 \RR_{\enfM} M_2$ if either:
  \begin{itemize}
   \item both $M_1,M_2$ diverge;
   \item $M_1 \Converge E_1[x V_1]$ and $M_2 \Converge E_2[x V_2]$
   for some $x$, values $V_1,V_2$, and evaluation contexts $E_1,E_2$ with
   $V_1 \RR_{\enfV} V_2$ and $E_1 \RR_{\enfK} E_2$;
   \item $M_1 \Converge V_1$ and $M_2 \Converge V_2$ for some values $V_1,V_2$
   with $V_1 \RR_{\enfV} V_2$.
  \end{itemize}
  \item $V_1 \RR_{\enfV} V_2$ if $V_1 x \RR_{\enfM} V_2 x$ for some fresh $x$;
  \item $E_1 \RR_{\enfK} E_2$ if $E_1[x] \RR_{\enfM} E_2[x]$ for some fresh $x$.
 \end{itemize}
 \noindent
 The largest \emph{enf-bisimulation} is called \emph{enf-bisimilarity}.
\end{definition} 